% SOURCE_AUTHORITY: sga5_fr_workpass.tex, lines 306--325 (LF-normalized unit).
% SOURCE_SCAN: SGA5 (1).pdf, Exposé I opening pages; scan page mapping pending frozen unit manifest.
\begin{center}
{\Large \textbf{SGA 5 --- Exposição I}}\\[0.4em]
{\large \textbf{COMPLEXOS DUALIZANTES}}\\[0.3em]
por A.~Grothendieck\\
(redação de L.~Illusie)
\end{center}

\section*{Introdução}

Esta exposição é dedicada ao teorema de bidualidade em cohomologia étale. A teoria desenvolvida aqui e em SGA~A~XVIII (SGA~A = SGA~4), para a topologia étale e os feixes construtíveis de torção sobre os esquemas, é formalmente muito semelhante à desenvolvida em HARTSHORNE~[H] para a topologia de Zariski e os feixes coerentes: esta última serviu-lhe, em grandíssima medida, de modelo. No entanto, enquanto no caso dos coeficientes ``contínuos'' a existência de complexos dualizantes não apresenta dificuldades, no caso dos coeficientes ``discretos'' ela é muito mais delicada de demonstrar e parece requerer essencialmente a resolução de singularidades. Veja-se, contudo, o teorema de bidualidade de Deligne (SGA~$4^{1/2}$~Teorema de finitude~4.3.).

O parágrafo~1 contém a definição dos complexos dualizantes e as propriedades formais que dela decorrem, em particular a interpretação dos funtores dualizantes como ``intercambiadores de funtores''.

No parágrafo~2 demonstra-se que, sobre um esquema localmente noetheriano e conexo, um complexo dualizante é determinado de maneira única, a menos de uma translação de graus e do produto tensorial com um feixe inversível.

No parágrafo~3 chega-se ao teorema central da exposição (3.4.1), que afirma que, sobre um ``bom'' esquema regular $X$, por exemplo um esquema excelente, noetheriano e regular de característica zero\footnote{Esta última restrição será desnecessária uma vez que se disponha, para esses esquemas, da resolução de singularidades ao estilo de Hironaka e do ``teorema de pureza''.}, o complexo $(\Z/n\Z)_X$ é dualizante.

O parágrafo~4 apresenta as aplicações da teoria ao teorema de dualidade local. Este exprime que, sobre um ``bom'' esquema $X$ provido de um ponto fechado $x$ cujo corpo residual é separavelmente fechado, e para um feixe construtível $F$, há um emparelhamento perfeito entre $\underline{\H}^i(F')_x$ e $\underline{\H}^{-i}_x(F)$, onde $F'$ é o dual de $F$, isto é, $F' = \R\,\underline{\Hom}(F,K_X)$, e $K_X$ é um complexo dualizante.

Por último, no parágrafo~5 demonstra-se diretamente que, sobre um esquema regular de dimensão~1, o complexo $(\Z/n\Z)_X$ é dualizante ($n$ primo com as características residuais).
